\documentclass[11pt]{article}
\usepackage[T1]{fontenc}
\usepackage{lmodern}
\usepackage{microtype}
\usepackage[margin=1in]{geometry}
\usepackage{amsmath,amssymb,amsthm,mathtools}
\usepackage{algorithm}
\usepackage{algpseudocode}
\usepackage{booktabs}
\usepackage{tabularx}
\usepackage[numbers,sort&compress]{natbib}
\usepackage[hidelinks]{hyperref}
\usepackage[capitalise,noabbrev]{cleveref}
\setlength{\emergencystretch}{3em}

\theoremstyle{plain}
\newtheorem{theorem}{Theorem}[section]
\theoremstyle{definition}
\newtheorem{definition}[theorem]{Definition}

\newcommand{\lean}[1]{\path{#1}}
\newcommand{\N}{\mathbb{N}}
\newcommand{\alive}{\mathrm{alive}}
\newcommand{\gone}{\mathrm{gone}}
\newcommand{\come}{\mathrm{come}}
\newcommand{\held}{\mathrm{held}}
\newcommand{\born}{\mathrm{born}}
\newcommand{\died}{\mathrm{died}}

\title{The Algebra of Cuts:\\
Positions, Contents and Their Differences over an Append-Only Log}
\author{Alberto Serrano-Calva}
\date{November 13, 2026}

\begin{document}
\maketitle

\begin{abstract}
A personal knowledge store records every change to its catalog in an append-only log, one chain a
provenance class, so a cut, one position a chain, names the store as of a moment. Its first read
of every version a cut holds built a row in Python a head, 1.7 to 1.8~s for the 400{,}274 heads
of a corpus chain whose walk in SQL takes 66~ms. This paper treats the cut as a lazy value with an
algebra instead. On positions, cuts form a distributive lattice, as the consistent cuts of a
distributed computation do, and the difference of two cuts is the events between them. On
contents, a version is one row of the log, alive on one half-open interval of positions, so the
union, the intersection and the difference of what cuts hold are predicates on a row's position
and its closing stamp: they compose without a read and run as one statement, and the signed
difference of two cuts is two range reads whose cost grows with what changed. The union on
contents is not the union on positions: it holds both versions of an edited note, and the two
agree exactly on rows that are never closed. Seven laws are proved in Lean~4, 110 theorems with
six controls among them, and a property test checks the code against a fold of the log. On a
copy of the corpus catalog, every head counts in 66~ms and reads as three columns in 0.22~s, and
the signed difference over a thousand positions is counted in 2.5~ms and read and sorted in 2.7
to 2.9~ms, against 3.6 to 3.8~s for two full reads and a set difference.
\end{abstract}

\section{Introduction}

The system is a personal knowledge store \cite{serranocalva2026palace}: notes, dialogue, papers and
code land through one writer into a dozen stores, and a sentence of a query algebra reads them.
Its catalog maps each source path to the content it holds now and is rewritten in place. A
landing log, built as the first stage of a design for reading the store as of a past cut, keeps
the catalog's past: every write appends a row to a hash-chained chain of its provenance class, in
the transaction of the catalog's own row, and a row stores the position of its path's next event
when that event lands, its closing stamp. A cut is one position a chain, about a dozen integers
for an instance.

That stage's read of every version a cut holds returned a list of row objects: 3.5 to 3.9~$\mu$s a
head beyond the walk, so the read passed a tenth of a second near 27{,}000 heads. Before a later
stage builds the store's graph at a cut on this read, its shape had to change. The design takes
five requirements as given. The value a cut names is the store's graph at its positions, evaluated
lazily, so that holding a cut reads nothing. Operators take one cut with another, a difference, a
union and an intersection among them, and operations of any kind apply to one cut or to a family of
cuts. Each of the two set operations has two forms, one of positions and one of contents, and both
are kept. An operation given no cut works on the current one, the way a stack operation takes the
value on top. And editing a node gives it a new version while it stays the same node, so that a cut
holds each node in the version alive there.

The contributions:
\begin{enumerate}
\item A design (\cref{sec:design}): the cut as an immutable frontier; operators typed so that every
  result that is a cut, a set or a difference is an operand again; four forms every operation is
  written on; and a default that binds the head once an evaluation.
\item The interval structure (\cref{sec:laws}): a version alive on one half-open interval makes
  every set operation a predicate on two columns, both differences of two cuts range reads, and
  the composition of differences exact through a cut between the two.
\item Seven laws, SNP-11 to SNP-17, proved in Lean~4 over a model of the log, with controls that
  show where each fails, and attacked by randomized instances.
\item A build over SQLite (\cref{sec:build}): statements bounded by the intervals so that the
  planner walks what changed, columns read as one buffer each, and a tool of the store's protocol
  for agents.
\item Measurements on a copy of a 400{,}274-item catalog (\cref{sec:eval}).
\end{enumerate}

\section{Background and related work}

\paragraph{Cuts.} A cut of a distributed computation is a set of events closed under the order of
each process; Chandy and Lamport recorded consistent global states with markers
\cite{chandy1985snapshots}, and Mattern showed that the consistent cuts form a lattice under union
and intersection, with vector time as its coordinates \cite{mattern1989virtual}. Down-sets of a
partial order form a distributive lattice \cite{davey2002lattices}. Each chain here is one append
order, so a cut is a down-set of a disjoint union of chains, the lattice is the product of the
chains' positions, and every frontier of chains that read no other chain is consistent. The log's
chains read no other chain but for a change of class, which writes a row in each of two chains in
one transaction (\cref{sec:limits}).

\paragraph{Temporal data.} Transaction-time databases keep each version of a record with the
interval of time it was current, and a survey compares the access methods for such data
\cite{salzberg1999access}; SQL:2011 stores a system-time period on each row of a versioned table
\cite{kulkarni2012sql2011}. The relational interval tree answers interval intersection queries
with the indexes a relational system already has \cite{kriegel2000intervals}, and a timeline index
serves temporal aggregation, joins and time travel in a column store \cite{kaufmann2013timeline}.
The landing log is a transaction-time table whose period ends at the closing stamp. Allen's
thirteen relations classify how two intervals stand \cite{allen1983intervals}; a cut asks only
whether a point lies in an interval, and that is what keeps every operator a range read.

\paragraph{Graphs over time.} DeltaGraph retrieves past snapshots of a graph from a hierarchy of
stored deltas \cite{khurana2013deltagraph}, a temporal graph algebra gives operators on evolving
graphs under point semantics \cite{moffitt2017tga}, and temporal networks are studied for the
structure of contacts in time \cite{holme2012temporal}. The cut here is a frontier over several
chains read under a grant, and its graph is a later stage; the quantities that stage updates over
a difference rest on standard tools: a low-rank correction of a factor \cite{woodbury1950inverting},
eigensolvers started from earlier vectors \cite{knyazev2001lobpcg}, and union-find for added edges
\cite{tarjan1975union}.

\paragraph{Tools.} The laws are proved in Lean~4 \cite{moura2021lean4} with mathlib
\cite{mathlib2020}; the code is tested with Hypothesis \cite{maciver2019hypothesis}.

\section{The model and the laws}\label{sec:laws}

\begin{definition}[Stores, cuts, reads]
A store is a family of lists indexed by chain keys $K$, and a cut is a function $c : K \to \N$.
The read at $c$ takes the first $c(k)$ events of each chain, the head is the cut of the lengths,
and the cuts inside a store are those at or under the head. The events between cuts $a$ and $b$
are, on each chain, those after position $a(k)$ and up to $b(k)$.
\end{definition}

\begin{definition}[Generations]
On a chain of keyed writes each write is a row $r$ with a key, a value, its position $\born(r)$,
and its closing stamp $\died(r)$, the position of the key's next write, or none. The row is
alive at $c$ when $\born(r) \le c$ and $\died(r)$ is none or above $c$. A tombstone and a purged
row carry no value: they close the row before them and are no version.
\end{definition}

\begin{definition}[Sets of rows]
A set is an expression: $\held(c)$, or the union, the intersection or the difference of two
sets. A row is in $\held(c)$ when it is alive at $c$, and in a combination by the Boolean
operation on the operands. The signed difference of two cuts is
$\gone(a,b) = \held(a) \setminus \held(b)$ and $\come(a,b) = \held(b) \setminus \held(a)$.
\end{definition}

The definitions are \lean{Gen}, \lean{Held} and \lean{Held.has} in
\lean{formal/Palace/SnapshotAlgebra/Defs.lean}, over the model of the reads in
\lean{formal/Palace/Snapshot/Defs.lean}. There is no complement, deliberately: every set lies
inside the contents of the cuts it names (\lean{Held.exists_atom_of_has}), and the complement, the
rows a cut does not hold, grows with every landing (\lean{Witness.complement_grows}).

\begin{theorem}[SNP-11, the lattice of cuts]
The cuts inside a store are closed under the pointwise meet and join, with the empty cut at the
bottom and the head at the top, and no cut leaves as the store grows. The read at $a \sqcup b$
is, chain by chain, the later of the two reads and at $a \sqcap b$ the earlier; the events
between $a$ and $b$ are those between $a \sqcap b$ and $b$, and those between $a$ and
$a \sqcup b$. (\lean{cutsIn}, \lean{readAt_sup}, \lean{readAt_inf}, \lean{readBetween_inf_left},
\lean{readBetween_sup_right}.)
\end{theorem}

\begin{theorem}[SNP-12, a row lives on an interval]
A row alive at two cuts is alive at every cut between them. For $a \le b$ on a chain,
\begin{align*}
r \in \come(a,b) &\iff a < \born(r) \le b \text{ and } \died(r) \notin [\born(r), b],\\
r \in \gone(a,b) &\iff \born(r) \le a < \died(r) \le b,\\
r \in \held(a) \cap \held(b) &\iff \born(r) \le a \text{ and } \died(r) \notin [\born(r), b],
\end{align*}
and a row is alive at every cut of a family exactly when it is alive at the least and the
greatest. On the generations of a chain a cut holds one row a key. (\lean{Gen.aliveAt_of_between},
\lean{Gen.aliveAt_and_not}, \lean{Gen.aliveAt_and_not'}, \lean{Gen.aliveAt_and},
\lean{Gen.forall_aliveAt_iff}, \lean{eq_of_key_eq_of_mem_contentAt}.)
\end{theorem}

The first two are what make a difference cheap: the rows come are a range of positions, read on
the primary key, and the rows gone a range of closing stamps, read on an index of them. Both cost
what changed, not what is held. The control is a value rather than a row: a path that returns to
bytes it held holds them at two cuts and not between (\lean{Witness.value_not_contiguous}), which
is why the element of a cut's contents is a row of the log and not a digest.

\begin{theorem}[SNP-13, the signed difference]
$\held(b) = (\held(a) \setminus \gone(a,b)) \cup \come(a,b)$. When $b$ lies between $a$ and $c$ on
every chain, in either direction,
\begin{align*}
\come(a,c) &= (\come(a,b) \setminus \gone(b,c)) \cup \come(b,c),\\
\gone(a,c) &= \gone(a,b) \cup (\gone(b,c) \setminus \come(a,b)),
\end{align*}
and a row that came and went on the way is on neither side. A sum over the rows at $b$ is the
sum at $a$ less the rows gone, with the rows come. Sorted by key, each row gone is removed or the
first of a replaced pair and not both, each row come is added or the second of one, and every key
of the difference is one the events between name. (\lean{Held.has_cut_eq_apply},
\lean{Held.has_come_trans}, \lean{Held.has_gone_trans}, \lean{Held.sum_rows_delta},
\lean{mem_gone_iff}, \lean{not_replacedOf_of_removedOf}, \lean{mem_come_iff},
\lean{key_mem_of_mem_delta}.)
\end{theorem}

Through a cut outside the two the composition fails (\lean{Witness.delta_through_outside}). The
region of identities the events between two cuts name holds more than the difference: a path
landed and retired between them is named by the events and holds a version at neither cut
(\lean{Witness.changed_not_delta}).

\begin{theorem}[SNP-14, closure]
Whether a row is in a set is a function of its position and its closing stamp alone; a set read
is the set operation on the reads; a row's own test commutes with every set; a set whose cuts
lie inside the chain is unchanged by its growth; and what every cut of a family holds is what
its least and its greatest both hold. A function of two reads, a measure at each cut of a family
and a fold over one are unchanged by growth. (\lean{Held.has_eq_of_interval},
\lean{Held.mem_rows_union}, \lean{Held.rows_filter}, \lean{Held.rows_gens_append},
\lean{Held.has_all_eq_kept}, \lean{combine_grow}, \lean{series_grow}, \lean{foldOver_grow}.)
\end{theorem}

\begin{theorem}[SNP-15, contents against positions]
What either of two cuts holds is what their join holds with the rows gone between their meet and
their join; what both hold is what their meet and their join both hold. For a row never closed,
the join's contents are the union of the contents and the meet's their intersection.
(\lean{Held.has_union_cut}, \lean{Held.has_inter_cut}, \lean{Gen.aliveAt_sup_of_open},
\lean{Gen.aliveAt_inf_of_open}.)
\end{theorem}

The control is one key written twice between two cuts: the union on contents holds two rows and
the contents of the join one (\lean{Witness.union_ne_join}). So the two meanings are two
operators, and the types tell them apart.

\begin{theorem}[SNP-16 and SNP-17, the grant and the erasure]
A set read on the chains a grant shows is the grant's restriction of the set read on the whole
store, reads only those chains and its cuts' positions on them, and names no other. The closing
stamps are a function of the keys alone, so a purged row that keeps its key moves no position and
no stamp; it leaves every set and changes no other row of any set. A retire is a removal in every
signed difference that spans it. (\lean{Held.across_restrict}, \lean{Held.across_congr},
\lean{gens_map}, \lean{Held.rows_gens_holes}, \lean{removed_of_retire}.)
\end{theorem}

Each proof is a case analysis on the Boolean values of aliveness at the cuts involved, or an
induction on the list of rows; the composition laws use contiguity to exclude the one case,
alive at $a$ and $c$ and not at $b$, that a cut between rules out. The library proves the seven
laws in 110 theorems, six of them the controls, with no axiom beyond the standard three; 80 of
the declarations are restated with their proofs left out and judged identical by a comparator,
and every proof is replayed by six independent kernel checkers. Seeds 529 to 566 attack the
statements with 500 random instances each: twenty-one attacks survive, one of them on the 295
instances its guard admits, and seventeen negative controls are refuted.

\section{The design}\label{sec:design}

\Cref{tab:objects} gives the objects and what each holds before anything is read. A cut is its
positions, compared and hashed as a value, so two equal cuts are one key. Every other object holds
cuts and how they are combined, and reading is a separate act: a count, columns, or the sorted
rows of a difference.

\begin{table}[t]
\centering\small
\begin{tabularx}{\linewidth}{@{}lXX@{}}
\toprule
object & holds, unread & read as \\
\midrule
a cut & one position a chain & nothing: it is the value \\
events between two cuts & the two cuts & a count; rows of every kind \\
a set of versions & an expression over cuts & a count; columns \\
a signed difference & two sets, gone and come & counts of removed, added, replaced; sorted rows \\
a family & a list of cuts & by operation \\
a number, a series, a table & nothing: a result & itself; no operand of a set operator \\
\bottomrule
\end{tabularx}
\caption{The objects of the algebra. Join and meet take two cuts and return a cut; union,
intersection, difference and symmetric difference take two sets and return a set, a cut given
there being the set it holds.}\label{tab:objects}
\end{table}

\paragraph{The interface.} Every operation is written against four forms: \emph{apply}, a
function of one cut's read; \emph{combine}, a function of two; \emph{map}, a function at each
cut of a family; and \emph{fold}, an accumulation over a family. A read hands an operator the cut,
the versions it holds and the events up to it, all unread. The named operations are instances:
what a cut holds is an apply, the events between and the signed difference are combines, what
persisted through a family is a fold that carries the family's meet and join, and the count at
each cut walked by the differences is a fold that carries the last read and its count. Each form
is a function of the prefixes, so by SNP-14 each is unchanged by growth, and any operation that
is none of the named ones is admitted by writing it against the same four.

\paragraph{The default.} An operation that names no cut reads the head, and the head is read once,
when the evaluation starts: every part of one answer reads one cut, and the answer names its
positions. Asked again with those positions, the question gives the same answer whatever has
landed since, up to what a later purge erased. The stack the requirement compares the default to
is a comparison only; no stack of cuts and no state is kept between two evaluations.

\paragraph{Scope and erasure.} The grant names chains, so every statement names the chains of the
grant's classes and no other; a cut naming another chain is refused in the words used for a chain
that does not exist. A purged version is a row without its salt, which the version's own test
leaves out of every set; the answer counts such rows in the ranges it read.

\section{The build}\label{sec:build}

The engine is one module over the log's tables in the catalog's SQLite file. An evaluation reads
the chains of the asker's classes and their heads in one statement; cuts are restricted to those
chains, and a position past a chain's head is refused.

\paragraph{From an expression to a statement.} For each chain the expression compiles to two
things (\cref{alg:compile}): the predicate, a Boolean combination of
\texttt{seq <= c AND (closed\_at IS NULL OR closed\_at > c)} for each cut's position $c$; and the
bounds the intervals imply, a range of positions and a range of closing stamps that every row of
the set lies in. The bounds add nothing to the predicate's meaning, but they are ranges an index
can take: for $\come(a,b)$ the positions $(a, b]$, on the primary key; for $\gone(a,b)$ the closing
stamps $(a, b]$, on the index of closing stamps. A chain whose bounds are empty is skipped, so a
difference between two cuts that agree on a chain reads nothing there. A set of any depth is then
one statement, and combining two sets runs none.

\begin{algorithm}[t]
\caption{The bounds of a set on one chain: positions in $(\ell, h]$, closing stamps in $(d, u]$,
$u = \top$ for an open row.}\label{alg:compile}
\begin{algorithmic}[1]
\Function{Bounds}{$e$}
  \If{$e = \held(c)$} \Return $(0, c, c, \top)$ if $c > 0$, else empty \EndIf
  \State $P \gets \Call{Bounds}{e.p}$
  \If{$e$ is a union} \Return the hull of $P$ and \Call{Bounds}{$e.q$} \EndIf
  \If{$e$ is an intersection} \Return the meet of $P$ and \Call{Bounds}{$e.q$}, or empty \EndIf
  \For{each cut $c > 0$ that $e.q$ is, or is a union of}\Comment{$e = e.p \setminus e.q$}
    \If{$d_P \ge c$} $\ell_P \gets \max(\ell_P, c)$\Comment{alive past $c$, not at it: stored after it}\EndIf
    \If{$h_P \le c$} $u_P \gets \min(u_P, c)$\Comment{stored by $c$, not alive at it: closed by it}\EndIf
  \EndFor
  \State \Return $P$, or empty if a range is empty
\EndFunction
\end{algorithmic}
\end{algorithm}

\paragraph{Columns.} A count is one statement of two aggregates, the versions and the holes, and
builds no row. Rows are read as columns: each column is aggregated in SQL into one string and
decoded into one buffer, an integer column into an array of 64-bit integers, the path's key into
an array of 32-byte values, and text into one run of bytes with an array of offsets, decoded a
value at a time when a listing asks for it. Reading rows makes fewer than one Python object a
row.

\paragraph{The signed difference.} Its counts are one statement: the two sides, materialized,
joined on the chain and the path's key, so a path replaced is counted once and removed and added
are the rest. Read, both sides are sorted in one sort of their 40-byte keys, the chain and the
path's key; since a cut holds one row a path on a chain, a key on both sides is two neighbours,
the row gone first. Composition takes the formulas of SNP-13 and is refused unless the shared cut
lies between the outer two on every chain.

\paragraph{Families.} What persisted reads two positions a chain whatever the family's size. The
pairwise table reads once what some cut of the family holds, as positions and closing stamps,
builds the matrix $M$ with $M_{ri} = 1$ when cut $i$ holds row $r$, and returns
$M^{\top}(1 - M)$, whose entry $(i,j)$ counts what $i$ holds and $j$ does not.

\paragraph{The tool.} One tool of the store's protocol for agents, \texttt{cut}, takes an operator
and its operands: each a cut's name (\texttt{head}, \texttt{bottom}, or positions written
\texttt{land.<chain>@<seq>}) or a call written as an object, whose result is the operand. So a
call is one expression: it reads the heads when it starts, a chain's horizon once for a position
named under the chain's head, and then what its outermost operator needs. The server opens the
catalog's file to only read, and a call's sets name at most 256 cuts, counted each time their
statement names one. A set given where a cut is expected is refused by its type, and the refusal
names what was given and what is taken. A
position under a chain's horizon, inside the block that opens a chain and is not history, is read
at the chain's empty prefix, and the answer's note says so. Every answer names the cuts it read
by their positions and says that it is of versions only: edges and graph mathematics at a cut
arrive with a later stage.

The read the contents replace is removed with nothing kept beside it; the measurements keep its
statement as their baseline. The tests check every law against a plain fold of the log kept in
the test code, over Hypothesis's random sequences of catalog writes: records, retires, renames,
relabels, re-keys, changes of class and purges, at every position and every pair of positions of
each chain and at random frontiers over several chains, ordered and not.

\section{Evaluation}\label{sec:eval}

The deterministic tier of the repository's tests passes at 3{,}757 tests with the algebra's 36.
The measurements are of the code as built, on two copy-on-write clones of a corpus instance's
catalog: 400{,}274 items of one class, so one chain whose genesis ends at position 400{,}276, then
200{,}000 writes after it (100{,}000 new versions of held paths, 50{,}000 retires, 50{,}000 new
paths, shuffled), the chain ending at 600{,}276 with 400{,}274 heads. One machine, an Apple M2
Pro with 32~GB, Python 3.15.2, SQLite 3.54.1 and NumPy 2.5.4, with other work running at a load
average between 2.3 and 4.0. The figures are medians of five repetitions, the two runs given as a
range (\cref{tab:heads,tab:delta}).

\begin{table}[t]
\centering\small
\begin{tabular}{@{}lrr@{}}
\toprule
every head of the chain, 400{,}274 & at the genesis's head & after the history \\
\midrule
a row in Python a head (the read replaced) & 1.74--1.79~s & 1.82--1.86~s \\
counted & 66~ms & 90--95~ms \\
three columns: position, key, digest & 219--224~ms & 236--241~ms \\
eight columns & 433--437~ms & 440--442~ms \\
\bottomrule
\end{tabular}
\caption{Every head of the corpus chain. The columns equal the rows in both runs.}\label{tab:heads}
\end{table}

\begin{table}[t]
\centering\small
\begin{tabular}{@{}rrrrrr@{}}
\toprule
positions & removed, added, replaced & counts & sides as columns & sorted & two reads \\
\midrule
1 & 1, 0, 0 & 0.02 & 0.02, 0.02 & 0.08 & 3{,}670--3{,}728 \\
1{,}000 & 248, 224, 528 & 2.5 & 0.4, 1.9--2.0 & 2.7--2.9 & 3{,}684--3{,}742 \\
100{,}000 & 24{,}995, 25{,}010, 49{,}995 & 279--305 & 42, 218--238 & 318--362 & 3{,}655--3{,}787 \\
200{,}000 & 50{,}000, 50{,}000, 100{,}000 & 560--604 & 83--90, 448--463 & 657--682 & 3{,}634--3{,}812 \\
600{,}276 & 0, 400{,}274, 0 & 144--147 & 238--260, 0.01 & 398--448 & 1{,}783--1{,}879 \\
\bottomrule
\end{tabular}
\caption{The signed difference ending at the head, over spans of positions, in milliseconds: its
counts, its two sides as columns (come, then gone), its rows sorted, and two reads of every head
with a set difference. Every side equalled the set difference.}\label{tab:delta}
\end{table}

\paragraph{Every head.} Counted, every head of the chain takes 66~ms, 0.17~$\mu$s an event walked
and no row built, 27 times faster than the rows it replaces, so a count passes a tenth of a second
up to about 600{,}000 events. Three columns take 0.22~s, eight times faster than the rows and over
that bar: beyond the walk a head costs 0.39~$\mu$s as three columns, against 4.2~$\mu$s as a row.
Three columns pass the bar near 180{,}000 heads on a chain whose events are its heads. That is the
size at which the design keeps heads at recorded cuts or walks the differences from a first read,
and the corpus chain is over it.

\paragraph{The difference.} Counted, it costs what changed: 0.02~ms for one position, 2.5~ms for
a thousand (2.7 to 2.9~ms read and sorted), 0.28 to 0.31~s for 100{,}000 and 0.56 to 0.60~s for
the whole history, against 3.6 to 3.8~s for two reads of every head and a set difference at every
span that ends at the head from inside the history, and 1.8 to 1.9~s from the bottom, where one
of the two reads is empty. The planner took the primary key for the rows come and the index of
closing stamps for the rows gone, unprompted. The
rows gone cost 2.9 to 3.2~$\mu$s a row against 0.57 for the rows come, since each is found on the
index and then read from the table by its position, where the rows come are one run of the
primary key.

\paragraph{A family.} For ten cuts spread over the history, the counts walked by the differences
took 0.43 to 0.49~s against 0.79 to 0.83~s for ten counts, and equalled them; what persisted
through all took 64~ms, one walk to the least position; what existed in any 99~ms; and the
pairwise table 0.32~s, equal to the hundred differences' sizes.

\section{Limitations}\label{sec:limits}

The model states the laws over lists, and the statements the code runs are not in it: that each
runs as its predicate is the property test's, not a theorem's. Every head as columns is over the
bar at the corpus's size, so a reader of every head of such a chain reads it once and walks the
differences. The rows gone are read through an index and the table, five times the cost a row of
the rows come. The measurements are of one chain of one class, so a difference between two cuts
neither of which is under the other was tested and not measured, and the history purged nothing.
A position named under a chain's head is checked against the chain's horizon by a walk back to
the horizon's row, once a call for each chain, since no index finds it. A change of class writes
a row in each of two chains in one transaction, and a frontier can hold one without the other,
so that the path is in both classes or in neither; such a frontier is read as asked and the tool
says so, and whether the pair becomes a reference the resolver follows is left to the next
stage. And this is not yet the graph: the edges by species and the graph's quantities at a cut
are the next stage's.

\section{Future work}

The next stages add the store's other chains (the chat log, the attestation ledger, the commit
order) to the same cut value, then build the graph at a cut on this algebra: the nodes as the
versions a cut holds, each handed to an operator as its parts, the edges by species and their
differences between two cuts, and the quantities updated over a difference where that is cheaper.
On synthetic nearest-neighbour graphs the design measured where each update pays: a low-rank
correction of a factor up to about a hundred changed edges, eigenvectors carried from the earlier
cut with its factor as the preconditioner at every edit measured, and union-find for added edges
up to a few hundred to a few thousand. A later stage writes the algebra's regions into the query
language, and whether an agent outside the owner's grant may call the tool is left to the owner.

\section{Conclusion}

A cut is a handful of positions. Reading what it holds as rows built in Python cost a row object
a head; reading it as a predicate on two columns costs a walk, and the difference of two cuts
costs what changed between them. The two meanings of union and intersection differ by a theorem
and by their types, the seven laws the code relies on are machine-checked, and the code is tested
against them on random histories.

\bibliographystyle{plainnat}
\bibliography{refs}
\end{document}
